% =====================================================================
%  Aufgabe 2
% =====================================================================
\chapter{Formative Evaluation}
\label{chap:formative_evalutation}

\section{Einleitung}
\label{sec:pt-einleitung2}

Mit der formativen Evaluation wollen wir herausfinden, wie wir tick.it weiterentwickeln können, um die Anwendung auf die Bedürfnisse von Studierenden anzupassen. Wir verfolgen bei der Evaluation die folgenden drei Forschungsfragen:

\begin{enumerate}
    \item[\textbf{FF1}] Wie organisieren Studierende ihre Aufgaben mit digitalen Tools und welche Schwierigkeiten erleben sie dabei?
    \item[\textbf{FF2}] Wie erleben Studierende die Bearbeitung eines einfachen Organisationsszenarios in tick.it und welche Stärken und Probleme zeigen sich dabei?
    \item[\textbf{FF3}] Welche Anforderungen stellen Studierende an ein kollaboratives Tool zur Organisation von Gruppenarbeit?
\end{enumerate}

Unsere primäre Zielgruppe sind Studierende, die ihren Studienalltag organisieren möchten. Aus diesem Grund setzen wir für die Teilnahme voraus, dass die Personen aktuell in einem Studiengang eingeschrieben sind. Zusätzlich erfassen wir zu Beginn des Interviews vier Eigenschaften, um die Aussagen später einordnen zu können. Das Fachsemester gibt Aufschluss darüber, wie viel Erfahrung die Teilnehmenden mit der Organisation ihres Studiums im Allgemeinen haben. Die bisherige Nutzung digitaler Organisationstools ermöglicht es, ihre Erfahrung mit solchen Tools einzuordnen. Ob im Semester mehrere Aufgaben parallel anfallen, ist relevant, da die Bearbeitung mehrerer Aufgaben ein Kernszenario eines Kanban-Boards darstellt. Zudem erfassen wir, wie häufig die Teilnehmenden in Gruppen arbeiten, um ihre Aussagen zu den Anforderungen an ein kollaboratives Organisationstool einordnen zu können.

Mit dem Ziel die Bedürfnisse der Nutzenden tiefer zu verstehen, führen wir eine thematische Analyse der Interviews durch. Da die Interviews qualitative Daten liefern, die Perspektiven und Erfahrungen der Teilnehmenden abbilden, eignet sich diese Methode besonders gut. Sie ist zudem offen genug, um auch unerwartete Themen zu erfassen, die wir mit einem vorab festgelegten Kategoriensystem übersehen könnten.


\section{Vorgehen}
\label{sec:pt-vorgehen}

NOTIZEN:
Zur Beantwortung der Forschungsfragen führen wir halbstrukturierte Interviews durch. Sie eignen sich, um Erfahrungen, Schwierigkeiten und Bedürfnisse der Teilnehmenden offen zu erfassen. Zusätzlich bearbeiten die Teilnehmenden eine konkrete Aufgabe in tick.it und sprechen dabei ihre Gedanken laut aus. So verbinden wir allgemeine Erfahrungen mit dem Eindruck aus der tatsächlichen Nutzung des Prototyps und können Stärken und Probleme direkt am Szenario beobachten.
Die Interviews wurden mit Whisper (OpenAI, 2025; Radford et al., 2023), Modell medium, lokal auf einem eigenen Rechner automatisch transkribiert, sodass keine Audiodaten an externe Server übertragen wurden. Anschließend wurden die Transkripte manuell anhand der Audioaufnahmen geprüft und korrigiert sowie nach den Transkriptionsregeln überarbeitet."

% ---------------------------------------------------------------------
\begin{beitraege}
  Text \cref{sec:pt-einleitung2}  & Tara Dethlefsen \\
  \midrule
  KI/Automatisierung: keine        & -- \\
\end{beitraege}
